\begin{pastexam}[label={px:df3-099-ds-01}]{地方三等資訊處理 99 年 第一題}
（申論；25 分）

解釋下列名詞並舉例說明：（每小題 5 分，共 25 分）
\begin{itemize}[leftmargin=2.2em]
\item[㈠] 演算法（algorithm）
\item[㈡] 時間複雜度（time complexity）
\item[㈢] 遞迴式的解決問題方法（recursive solution）
\item[㈣] 雙向佇列（Deque）
\item[㈤] 最小成本生成樹（minimum cost spanning tree）
\end{itemize}
\end{pastexam}

\begin{pastsoln}
五個名詞分屬資結書第~1~章（㈠㈡）、第~2~章與第~6.4~節（㈢）、第~6.2.6~節（㈣）、第~24~章（㈤），本書放在第~1~章。
名詞解釋題要寫「定義＋要點＋一個具體例子」。

\textbf{㈠ 演算法}。 解決某個問題的一串明確、有限的步驟：給定輸入，照步驟做，在有限步之後停下來並產生正確的輸出。
通常要求五個性質：\textbf{輸入}（零個或多個）、\textbf{輸出}（至少一個）、\textbf{明確性}（每一步意思清楚、沒有歧義）、
\textbf{有限性}（一定在有限步內結束）、\textbf{有效性}（每一步都是做得到的基本操作）。 例：求兩個正整數最大公因數的
輾轉相除法——當 $b \ne 0$ 時以 $(b, a \bmod b)$ 取代 $(a, b)$，$b = 0$ 時輸出 $a$；$\gcd(48, 18)$：$(48, 18) \to (18, 12)
\to (12, 6) \to (6, 0)$，輸出 $6$。 每一步 $b$ 都變小，所以一定會停。

\textbf{㈡ 時間複雜度}。 演算法執行所需的基本操作次數，表示成輸入大小 $n$ 的函數 $T(n)$，通常只看它的成長階，用
漸近符號 $O$、$\Omega$、$\Theta$ 表示（資結書第~1.5~節），並區分最壞、平均、最好情況。 它和機器速度無關，用來比較
演算法在 $n$ 變大時的表現。 例：在 $n$ 個數的陣列中循序搜尋 $x$，最壞情況要比較 $n$ 次（$x$ 在最後或不存在），時間
複雜度 $O(n)$；若陣列已排序改用二元搜尋，最壞只要約 $\lfloor \log_2 n \rfloor + 1$ 次，是 $O(\log n)$。 $n = 10^6$ 時前者
最多一百萬次、後者最多 $20$ 次。

\textbf{㈢ 遞迴式的解決問題方法}。 把問題化成\textbf{同類型但規模較小}的子問題，呼叫自己去解子問題，再把子問題的答案
組合成原問題的答案；並且有不再呼叫自己、直接得到答案的\textbf{終止條件}（base case）。 每次遞迴都要往終止條件靠近，
否則不會停。 執行時系統用堆疊保存每一層呼叫的參數和返回位置（資結書第~6.4~節）。 例：階乘
$n! = n \times (n-1)!$、$0! = 1$，寫成 \proc{Fact}$(n)$：若 $n = 0$ 傳回 $1$，否則傳回 $n \cdot$\proc{Fact}$(n-1)$；
\proc{Fact}$(3) = 3 \cdot$\proc{Fact}$(2) = 3 \cdot 2 \cdot$\proc{Fact}$(1) = 3 \cdot 2 \cdot 1 \cdot$\proc{Fact}$(0) = 6$。
河內塔、合併排序、二元樹走訪也是典型的遞迴解法。

\textbf{㈣ 雙向佇列（deque，double-ended queue）}。 兩端都可以加入和刪除元素的線性串列，支援四個操作：前端加入、
後端加入、前端刪除、後端刪除，都是 $O(1)$（資結書第~6.2.6~節）。 只用一端加入和刪除時就是堆疊；後端加入、前端
刪除時就是佇列，所以 deque 同時涵蓋兩者。 常用環狀陣列或雙向鏈結串列實作。 例：空的 deque 依序做「後端加入 $1$、
後端加入 $2$、前端加入 $0$」得 $\langle 0, 1, 2 \rangle$；再「後端刪除」得到 $2$、「前端刪除」得到 $0$，剩下
$\langle 1 \rangle$。 應用如瀏覽器的「上一頁／下一頁」記錄、滑動視窗最大值。

\textbf{㈤ 最小成本生成樹}。 連通、邊有權重（成本）的無向圖 $G = (V, E)$ 中，\textbf{生成樹}是包含全部 $|V|$ 個頂點、
恰好 $|V| - 1$ 條邊、沒有環的子圖；所有生成樹中\textbf{邊的成本總和最小}的就是最小成本生成樹（資結書第~24~章）。
常用 Kruskal（由小到大挑不成環的邊）或 Prim（從一點出發，每次加入連到樹外最便宜的邊）演算法求得。 例：四個城市
$a, b, c, d$ 之間可以鋪設的線路成本是 $ab = 1$、$bc = 2$、$ac = 3$、$cd = 4$、$bd = 5$，要讓四個城市都連通又最省：
Kruskal 依序選 $ab$（1）、$bc$（2），$ac$ 會和 $ab, bc$ 成環跳過，再選 $cd$（4），共 $3$ 條邊、總成本 $7$，就是
最小成本生成樹。
\end{pastsoln}
